% LSPV-3 REPORT::ROUNDABSORB1::2026-10-03T23:31:00-07:00
\documentclass[11pt]{article}
\usepackage[margin=.85in]{geometry}
\usepackage{amsmath,amssymb}
\begin{document}
\section*{LSPV-3 ROUNDABSORB1: polynomial-round local PH interaction is absorbed under P=NP}
\textbf{STATUS: OPEN.} P versus NP remains open. The theorem is \textbf{CONDITIONAL} on P=NP and awaits independent red-team review. It proves no SAT lower bound and no nonrelativizing or nonalgebrizing hinge.

\section*{CORPUS REVIEWED}
This run surveyed the repository tree and read root main.tex, AGENTS.md, CODEX.md, the pnp-research skill, current tasks, materially relevant LSPV reports, FULLSKOLEM1, PHABSORB1, the reconstruction ledger, HEC/CRS-PPC/CCSPC, and governance notes. The predecessor LSPV-2 CHAINEX1 body remains unverified in GitHub; LSPV-2 was instructed to retry only its exact preserved body if available and never reconstruct it from summary.

Primary literature checked includes Akshay et al. (SAT 2026) on Skolem/Herbrand extraction from a specified QBF proof system, Boehm--Beyersdorff (SAT 2026) on CAQE strategy extraction, and Gajulapalli--Ghentiyala--Li--Saraogi (ECCC TR25-121, 2025) on downward self-reducibility in the total function polynomial hierarchy. These are model-specific search/strategy results and do not provide the arbitrary-static-SAT bridge.

\section*{Exact interactive selector theorem}
Fix a SAME\_REALIZATION input $X$. Let $r(|X|)$ be polynomial. At round $i\le r(|X|)$ let the transcript prefix be $t_i=(u_1,w_1,\ldots,u_i)$, with polynomially bounded challenges and responses. Let $L_i(X,t_i)$ be a uniformly fixed-level PH legality predicate and $V_i(X,t_i,w_i)$ a uniformly fixed-level PH current-response predicate. Assume local totality on every reachable legal prefix:
\[
\forall i\le r(|X|)\ \forall t_i\quad L_i(X,t_i)\Rightarrow\exists w_i\ V_i(X,t_i,w_i).
\]
\paragraph{Theorem ROUNDABSORB1.} \textbf{CONDITIONAL on P=NP.} Under P=NP there is one deterministic polynomial-time online responder $F$ such that on every reachable legal prefix, $w_i=F(X,i,t_i)$ satisfies $V_i(X,t_i,w_i)$. The entire polynomial-round interaction uses polynomial total time.

\paragraph{Proof.} Under P=NP, $PH=P$. For a prefix $p$ of $w_i$, define $E_i(X,t_i,p)\equiv\exists s\ V_i(X,t_i,ps)$. Uniform fixed-level PH closure gives $E_i\in PH=P$. Bitwise prefix self-reduction constructs $w_i$; local totality guarantees continuation. Polynomial response length gives polynomial work per round, and polynomially many rounds preserve polynomial total work. \hfill$\square$

\section*{Boundary: local totality versus global continuation}
The theorem does not say every polynomial-round game has a polynomial-time winning strategy. If a move is acceptable only when it has a winning continuation through all future adversarial challenges, the continuation predicate may have quantifier depth growing with the polynomial horizon and can leave fixed-level PH for PSPACE-style alternating computation. ROUNDABSORB1 applies only when current acceptance is locally total at a uniformly fixed PH level, or when global continuation viability is separately proved to remain there.
\[
\boxed{\text{Adaptivity alone is not an escape; growing alternation depth or a non-PH forced resource might be.}}
\]

\section*{HEC / CRS-PPC / CCSPC consequence}
Fresh challenges, stopping-time randomness, delayed audits, and polynomially many adaptive epochs do not by themselves evade PHABSORB1. If after each challenge the demanded response is any member of a polynomially balanced, locally total, fixed-level PH relation, ROUNDABSORB1 gives an online polynomial responder under P=NP. No prediction of future randomness is needed.

A surviving causal acquisition theorem must show at least one of: current response validity is not fixed-level PH; local totality fails and only global continuation winning suffices; continuation viability escapes fixed-level PH; or the lower-bounded resource is independently forced on every static SAT decider. The last remains the TC4/SAME\_REALIZATION bridge.

\section*{Quantifier and accounting audit}
The theorem never commutes $\forall u_i\exists w_i$ and never precommits $w_i$ before $u_i$ is known. It does not compile the online process into one static object and therefore does not solve TASK-004 or TASK-005. All reconstruction charges remain separate: Q acquisition, normalization, census/output, amplification, indexing, compilation, route work, writing, and transport. TransportCert remains accounting only. ROUNDABSORB1 supplies no certified numerical exponent.

\section*{PQ and proof-complexity ledger}
Akshay et al. and Boehm--Beyersdorff are \textbf{KNOWN THEOREMS} in their declared proof-system models and only analogies for arbitrary-SAT acquisition. Construction-specific collision extraction remains an \textbf{ANALOGY/CANDIDATE}. Generic cryptographic binding remains \textbf{FALSIFIED/RETIRED} as an unconditional transfer. Gajulapalli et al. is a \textbf{KNOWN THEOREM} about downward self-reducible total search and reinforces that total search structure can yield upper bounds; it does not imply TC4.

\section*{Hostile controls and barriers}
DEAD, COPY, DELAY, INTERP, and RECOMP remain active. RECOMP is decisive when each round admits an independently searchable valid response. Horn, 2-SAT, affine/Tseitin, bounded treewidth, compact decision diagrams, structured circuits, memoization, and arithmetic representations remain controls. ROUNDABSORB1 relativizes: for any oracle $O$, $P^O=NP^O$ implies $PH^O=P^O$, and the same selector works relative to $O$. Algebraic-oracle augmentation supplies no lower-bound hinge.

\section*{DUPLICATION / CONTRADICTION AUDIT}
ROUNDABSORB1 extends PHABSORB1 from one challenge-response call to polynomially many adaptive calls under local totality. It does not duplicate FULLSKOLEM1, which handles the displayed fixed reconstruction prefix, and does not contradict the CRS-PPC Fresh Snapshot Lemma: freshness remains true, but the responder waits for each revealed challenge.

\section*{WHAT CHANGED RELATIVE TO ROOT MAIN}
Root main records FULLSKOLEM1 but not verified PHABSORB1. PHABSORB1 already closes total polynomially balanced fixed-level PH provenance search under P=NP. ROUNDABSORB1 adds that polynomially many adaptive rounds do not restore hardness when each round remains locally total and uniformly fixed-level PH.

\section*{ADOPTED / REJECTED / RETIRED PRIOR COMPONENTS}
\textbf{Adopted:} FULLSKOLEM1, verified PHABSORB1, SAME\_REALIZATION discipline, conservative exponent accounting, Fresh Snapshot at its declared scope, HEC/CRS-PPC/CCSPC as restricted models or candidate bridges.

\textbf{Rejected as sufficient hinges:} polynomially many adaptive challenges, fresh randomness alone, delayed audit alone, and polynomial response provenance when each round remains locally total fixed-level PH.

\textbf{Retained OPEN:} global continuation-winning predicates whose complexity grows with the horizon, universal static-to-online compilation, source-matched acquisition necessity, TC4, nonrelativization, and nonalgebrization.

\section*{STALE OR OVERSTATED CLAIMS FOUND}
Any statement that online freshness or polynomial-round adaptivity alone defeats PH semantic reacquisition is overstated. The missing theorem must identify a global continuation predicate outside fixed-level PH or a universally forced computation resource not reducible to response search.

\section*{One next target}
For the strongest CCSPC/CRS-PPC candidate, write the exact continuation predicate $W_i(X,t_i,w_i)$ and determine its uniform complexity as the remaining horizon grows. If it remains at a fixed PH level, ROUNDABSORB1 absorbs it under P=NP. If its alternation depth grows, quantify the resulting class and test whether the static-SAT compiler forces arbitrary SAT deciders to realize that global game rather than bypass it.

\section*{References}
S. Akshay et al., Proof Systems for QBF Synthesis: Extracting Skolem and Herbrand Functions, SAT 2026, LIPIcs 377, Article 3.

B. Boehm and O. Beyersdorff, Towards Understanding the Complexity of CAQE: A Proof-Theoretic Analysis of Its Core Procedure, SAT 2026, LIPIcs 377, Article 7.

K. Gajulapalli, S. Ghentiyala, Z. Li, and S. Saraogi, Downward self-reducibility in the total function polynomial hierarchy, ECCC TR25-121, 2025.

\section*{Final status}
ROUNDABSORB1 is \textbf{CONDITIONAL} on P=NP pending independent red-team review. SAME\_REALIZATION necessity, TC4, global continuation complexity, static-to-online compilation, source-matched no-sharing, and nonrelativizing/nonalgebrizing hinges remain \textbf{OPEN}. P versus NP remains \textbf{OPEN}.

% END_OF_VERIFIED_REPORT::LSPV-3::ROUNDABSORB1
\end{document}
